\chapter{Types, Constants, and Operations}
\label{ch:pebble-operations}

An opcode name is not an instruction specification.
The specification must also say which operands exist, what they mean, when they are read, and which result type is produced.
Otherwise a parser, verifier, and interpreter can each implement a plausible but incompatible version of \code{add}.

\section{A closed instruction set}
\label{sec:pebble-op-table}
Table~\ref{tab:pebble-ops} is the semantic core.
Input is allowed at most once, in entry, and has type \code{u8}.
Constants carry their result type explicitly.
Both phi inputs have the result's type, and their predecessor labels describe edges rather than operand evaluation order.

\begin{table}[htbp]
\caption{PebbleIR's live value-producing operations have complete operand contracts.}
\label{tab:pebble-ops}
\centering
\begin{tabular}{llll}
\toprule
Operation & Inputs & Result & Meaning \\
\midrule
input & none & u8 & Function argument \\
const & immediate & u8 or i1 & Typed literal \\
add & u8, u8 & u8 & Sum modulo 256 \\
mul & u8, u8 & u8 & Product modulo 256 \\
lt & u8, u8 & i1 & Unsigned strict comparison \\
phi & two same-type values & same type & Predecessor-selected value \\
\bottomrule
\end{tabular}
\end{table}

Terminators do not define values.
A return consumes one byte value; a jump names one successor; a branch consumes one boolean and names two distinct successors.
Edges into entry are prohibited.
The distinct-successor restriction keeps phi inputs keyed by predecessor blocks rather than by individual parallel edges.
Relaxing it would require a decision about edge identity throughout the representation, not only a change to the branch parser.

\section{Constants have types too}
\label{sec:pebble-constants}
The literal 1 can inhabit either type, but a \code{u8} constant 1 cannot directly supply an \code{i1} use.
For example, comparing byte constants 7 and 9 produces boolean true; folding that comparison must produce an \code{i1} constant, not a byte constant with the same bits.
The optimizer retains the original result type while replacing the opcode and clearing operands.

The implementation does not intern constants globally.
Two identical literals may occupy different slots and have different dominance regions.
Global interning would need either function-entry materialization or a separation between literal objects and instructions that materialize them.
Sharing the payload is a storage decision; sharing a dominating definition is a control-flow decision.

\begin{algorithm}[htbp]
\caption{Checking an operation against its semantic signature.}
\label{alg:pebble-signature}
\begin{pseudocode}
\Require a structurally bounded value record $v$
\Ensure compatible types and canonical fields, or rejection
\State dispatch on the validated opcode
\State require the documented result type and operand count
\State require zero in fields unused by that opcode
\State require each operand ID to name a live value
\State compare operand types with the signature
\State separately check where each operand is available
\end{pseudocode}
\end{algorithm}

Algorithm~\ref{alg:pebble-signature} separates type validity from dominance.
Both are necessary.
A correctly typed byte value from the left branch is not available on the right branch merely because its ID exists.
The verifier therefore validates all record headers before it follows operand IDs and before it computes availability queries over those records.

\section{Adding an opcode is a cross-client change}
\label{sec:pebble-op-extension}
Suppose the next client requests subtraction.
First define whether it wraps, traps, or has an undefined case.
For a byte-wrapping subtraction, specify the result as $(a+256-b)\bmod256$ and justify the host expression's range.
Then update the enum, textual name table, structural verifier, type rules, interpreter, constant folder, lowering, VM handler, and tests.
The current implementation deliberately has no subtraction; this procedure is an extension exercise, not a claim of an implemented opcode.

A generated opcode table could reduce duplicated signature metadata.
It would not generate a sound semantic implementation automatically.
In particular, a flag such as \enquote{pure} must describe behavior, not simply mirror a developer's wish to remove an instruction.

\section{Exercises}
\label{sec:pebble-four-exercises}
\begin{enumerate}
\item \label{ex:pebble-bool} Fold \code{lt 7,9} and explain which type and immediate the replacement needs. Why does constant byte 1 fail the branch contract?
\item \label{ex:pebble-sub} Work out wrapping subtraction for $(a,b)=(0,1)$ and $(255,255)$, and bound the proposed host intermediate $a+256-b$.
\end{enumerate}

\section{Further reading}
\label{sec:pebble-four-reading}
The C standard's rules for promotions and unsigned arithmetic explain why source and host types must be considered separately~\cite[Sections~6.3.1.1 and~6.2.5]{wg14-2011-c11}.
